% !PRINTED 39
\weberpage{39}
%<<CONT

\gesperrt{whole rational functions $F(x),\varphi(x)$ can be so determined, with $\varphi(x)$ of lower degree than $f(x)$, that the equation}
\begin{equation}
 F(x)f(x)+\varphi(x)\psi(x)=\chi(x)
\tag{26}
\end{equation}
\gesperrt{is identically satisfied.}

\begin{center}
{\large \S.\ 7.}

\vspace{0.5em}
{\large\gesperrt{Products of Linear Factors.}}
\end{center}

According to \S~1, by multiplication of whole functions we again obtain such functions of higher degree, and the degree of the product is determined by the sum of the degrees of the individual factors.

If therefore we multiply $n$ linear factors by one another, there arises a function of the $n$th degree, whose manner of formation we must examine somewhat more closely.

We shall take the linear factors in the simple form $x-\alpha_1$, $x-\alpha_2$, $\ldots$, $x-\alpha_n$, and put, since the coefficient of the highest power is (by \S~1) $=1$,
\begin{equation}
\begin{aligned}
 f(x)&=(x-\alpha_1)(x-\alpha_2)\cdots(x-\alpha_n)\\
 &=x^n+a_1x^{n-1}+a_2x^{n-2}+\cdots+a_n.
\end{aligned}
\tag{1}
\end{equation}
A slight consideration shows the following law of formation of the coefficients $a_1,a_2\ldots a_n$:

$-a_1$ is equal to the sum of the $\alpha$'s; $a_2$ is equal to the sum of the products of every two of the $\alpha$'s; $-a_3$ is the sum of the products of every three of the $\alpha$'s; in general $(-1)^\nu a_\nu$ is the sum of the products of every $\nu$ of the quantities $\alpha$, or, in formulas,
\begin{equation}
\begin{aligned}
-a_1&=\sum \alpha_1,\\
+a_2&=\sum \alpha_1\alpha_2,\\
&\cdots\\
(-1)^\nu a_\nu&=\sum \alpha_1\alpha_2\cdots\alpha_\nu,\\
&\cdots\\
(-1)^n a_n&=\alpha_1\alpha_2\cdots\alpha_n.
\end{aligned}
\tag{2}
\end{equation}
But to see still more clearly the correctness of this law of formation, one uses \gesperrt{complete induction}.

One first verifies the correctness in the first cases by actually carrying out the multiplication
%CONT>>
